\documentclass[12pt]{article}

\usepackage{sbc-template}
\usepackage[T1]{fontenc}
\usepackage[utf8]{inputenc}
\usepackage[brazilian]{babel}
\usepackage{graphicx,url}
\usepackage{float}
\usepackage{multirow}
\usepackage{booktabs}
\usepackage{amsmath}
\usepackage[none]{hyphenat}
\usepackage{tikz}
\usetikzlibrary{positioning,arrows.meta}

\sloppy
\hyphenpenalty=10000
\exhyphenpenalty=10000
\emergencystretch=3em
\setlength{\textfloatsep}{6pt}
\setlength{\floatsep}{6pt}
\setlength{\intextsep}{6pt}
\setlength{\abovecaptionskip}{4pt}
\setlength{\belowcaptionskip}{0pt}
\setlength{\parskip}{2pt}

\title{Segmentação Renal, Curadoria \textit{Human-in-the-Loop} e Quantificação da Ecogenicidade em Ultrassonografia: um Pipeline Rastreável para Caracterização Intrarrenal}

\address{
  CEFET/RJ -- Rio de Janeiro -- RJ -- Brasil\\
  $^{2}$ HUPE -- UERJ -- Rio de Janeiro -- RJ -- Brasil\\
\vspace{-0.7cm}
\email{jefferson.anjos.1@aluno.cefet-rj.br,}\\
\vspace{-0.7cm}  \email{\{amaro.lima,gabriel.araujo,felipe.henriques\}@cefet-rj.br},\\
\vspace{-0.7cm}  \email{\{jorge.henriques,norderval.norderval\}@uerj.br}
}

\begin{document}

\maketitle

\begin{abstract}
Renal ultrasound is inexpensive, non-ionizing and widely available, but computational analysis remains difficult because of speckle noise, acquisition variability, ambiguous anatomical boundaries and limited expert annotations. This work presents the evolution of a reproducible pipeline combining supervised kidney segmentation, external-data engineering, pseudo-label generation, model-consensus based prioritization, human-in-the-loop curation, multiclass intrarenal segmentation and quantitative echogenicity analysis. The kidneyUS dataset is adopted as the canonical supervised source. In a deduplicated test set of 70 images, U-Net achieved Dice 0.9290 and IoU 0.8722 for renal capsule segmentation. On 4,479 external MONAI images, U-Net generated 3,534 candidate renal pseudomasks. A multiclass U-Net achieved mean Dice 0.7594 for cortex, medulla and central echo complex on 50 manually annotated test images. A controlled medulla pseudo-label expansion increased the training set from 236 to 578 images but did not improve manual test Dice (0.7528 to 0.7523), reinforcing the need for human curation. Curated anatomical masks support an exploratory cortex--central echo complex echogenicity protocol measured on the original B-mode image. The workflow explicitly separates automation, curation and clinical interpretation and does not claim automatic diagnosis of chronic kidney disease or fibrosis.
\end{abstract}

\begin{resumo}
A ultrassonografia renal é acessível, não ionizante e amplamente disponível, mas sua análise computacional permanece difícil devido ao ruído \textit{speckle}, à variabilidade de aquisição, à ambiguidade dos limites anatômicos e à escassez de anotações especializadas. Este trabalho apresenta a evolução de um pipeline reprodutível que integra segmentação supervisionada do rim, engenharia de dados externos, geração de pseudorrótulos, priorização por consenso entre modelos, curadoria \textit{human-in-the-loop}, segmentação intrarrenal multiclasse e análise quantitativa de ecogenicidade. O kidneyUS é adotado como fonte supervisionada canônica. Em um teste deduplicado de 70 imagens, a U-Net obteve Dice 0,9290 e IoU 0,8722 na cápsula renal. Em 4.479 imagens externas do MONAI, foram produzidas 3.534 pseudomáscaras candidatas. Na etapa intrarrenal, a U-Net multiclasse alcançou Dice médio de 0,7594 em 50 imagens de teste manual. Uma expansão controlada de medula ampliou o treino de 236 para 578 imagens, mas não melhorou o Dice no teste manual (0,7528 para 0,7523), reforçando a necessidade da curadoria humana. Máscaras anatômicas revisadas sustentam um protocolo exploratório de ecogenicidade córtex--CEC medido na imagem B-mode original. O fluxo separa automação, curadoria e interpretação clínica e não afirma diagnóstico automático de doença renal crônica ou fibrose.
\end{resumo}

\section{Introdução}

A ultrassonografia renal combina baixo custo, portabilidade, ausência de radiação ionizante e aquisição em tempo real, mas sua interpretação depende de fatores como ruído \textit{speckle}, sombras acústicas, ganho, profundidade, equipamento e enquadramento. Essas características também tornam a análise computacional difícil: sem uma delimitação anatômica confiável, medidas de intensidade e textura podem incorporar fundo, textos da tela, tecidos adjacentes ou artefatos.

A segmentação externa do rim constitui, portanto, o primeiro nível do problema. A partir de uma máscara renal confiável, a análise pode ser restringida ao córtex, medula e complexo ecogênico central (CEC). O segundo desafio é a disponibilidade de rótulos. Revisões sobre segmentação médica semi-supervisionada mostram que a anotação pixel a pixel é cara e especializada, motivando pseudorrótulos, consistência entre modelos e seleção de casos\~\cite{Jiao2024SSLReview,Tajbakhsh2020Imperfect}. Entretanto, rótulos automáticos incorretos podem reforçar erros; trabalhos recentes enfatizam a confiabilidade e a correção explícita de pseudorrótulos\~\cite{Su2024ReliablePseudo,Huang2025SemiPLC}.

Este projeto surgiu como uma comparação de segmentadores renais, mas evoluiu para uma infraestrutura mais ampla: engenharia de dados, deduplicação, segmentação da cápsula, expansão por fontes externas, consenso entre arquiteturas, segmentação intrarrenal, curadoria \textit{human-in-the-loop}, inferência integrada e quantificação da ecogenicidade. A pergunta deixou de ser apenas qual modelo segmenta melhor o rim e passou a incluir como transformar predições em uma cadeia rastreável de imagem, anatomia, revisão humana e medida quantitativa.

O kidneyUS foi adotado como fonte supervisionada canônica por oferecer imagens B-mode e duas séries independentes de anotações para cápsula, córtex, medula e CEC\~\cite{Singla2023KidneyUS}. A etapa de ecogenicidade é motivada por estudos quantitativos que relacionam medidas relativas de brilho à função renal\~\cite{Manley2001,Beutler2024SRI,Zhao2025Echogenicity}. Ao mesmo tempo, correlações entre ecogenicidade e histopatologia são moderadas e inespecíficas\~\cite{Moghazi2005}; por isso, a presente pesquisa não trata ecogenicidade como diagnóstico de fibrose.

O objetivo geral é desenvolver um pipeline rastreável para segmentação e caracterização quantitativa da ultrassonografia renal combinando aprendizado profundo e curadoria humana. Os objetivos específicos são: estabelecer um benchmark deduplicado da cápsula; organizar fontes externas preservando proveniência; avaliar consenso entre arquiteturas; segmentar córtex, medula e CEC; estudar a propagação de erro na cascata; avaliar o impacto de pseudorrótulos; implementar curadoria auditável; e definir um protocolo exploratório de ecogenicidade córtex--CEC.

\section{Fundamentação e Trabalhos Relacionados}

\subsection{Bases abertas e ultrassonografia renal}

O kidneyUS foi proposto para reduzir a escassez de bases abertas em ultrassonografia renal, reunindo imagens provenientes de diferentes fabricantes e anotações independentes de especialistas\~\cite{Singla2023KidneyUS}. Essa característica permite que o conjunto seja usado tanto para benchmark de segmentação quanto para estudo de concordância entre anotadores. Como fonte externa não rotulada, este projeto utilizou o repositório clínico MONAI/NVIDIA, mantendo-o separado do benchmark manual.

\subsection{Segmentação renal}

A U-Net permanece uma referência em segmentação biomédica\~\cite{Ronneberger2015}; UNet++ amplia as conexões de salto\~\cite{Zhou2018UNetPP}; DeepLabV3 utiliza convoluções atrous para contexto multiescala\~\cite{Chen2018DeepLab}; e SegFormer introduz uma formulação baseada em transformadores\~\cite{Xie2021SegFormer}. Em ultrassom renal, Khan et al. destacam baixo contraste, ruído, variação de forma e artefatos como dificuldades centrais e mostram a relevância de pré e pós-processamento\~\cite{Khan2024KidneySeg}. Trabalhos recentes também utilizam o kidneyUS para segmentação multiclasse de estruturas internas\~\cite{Zhang2025RenalAttention}.

\subsection{Pseudorrótulos, consistência e semi-supervisão}

A pseudorrotulação utiliza previsões de um modelo como rótulos candidatos para dados não anotados\~\cite{Lee2013PseudoLabel}. Em segmentação médica, métodos recentes combinam pseudorrótulos com consistência, incerteza, múltiplos ramos ou correção de rótulos\~\cite{Jiao2024SSLReview,Su2024ReliablePseudo,Huang2025SemiPLC}. A metodologia deste trabalho adota uma postura conservadora: confiança, morfologia e consenso ajudam a ordenar a revisão, mas não transformam automaticamente uma máscara em referência manual.

\subsection{Human-in-the-loop e curadoria}

Budd et al. descrevem \textit{human-in-the-loop} como estratégia relevante em análise médica, incluindo seleção de casos, interação com previsões e refinamento iterativo\~\cite{Budd2021HITL}. Neste projeto, o revisor recebe uma proposta automática e pode aceitar, corrigir, rejeitar ou marcar uma estrutura como indisponível. A proposta original é preservada e cada edição fica associada ao revisor e ao tempo.

\subsection{Ecogenicidade e função renal}

Manley e O'Neill demonstraram que a ecogenicidade cortical pode ser quantificada de forma reprodutível em relação a uma referência, embora seja influenciada por condições de aquisição\~\cite{Manley2001}. Beutler et al. estudaram o índice esplenorrenal e encontraram associação com creatinina, ureia e eGFR\~\cite{Beutler2024SRI}. Zhao et al., em 339 participantes, relataram associação entre diferenças seio--córtex e função renal\~\cite{Zhao2025Echogenicity}. Esses resultados apoiam medidas relativas, mas não um valor absoluto universal.

\subsection{Fibrose e IFTA}

Moghazi et al. compararam ultrassom e biópsia em 207 pacientes e observaram que ecogenicidade teve as maiores correlações sonográficas com alterações histológicas, porém com magnitude moderada e sem fibrose intersticial como determinante independente no modelo multivariado\~\cite{Moghazi2005}. Athavale et al. desenvolveram um modelo para IFTA usando biópsia como referência\~\cite{Athavale2021}, e Qin et al. mostraram melhor desempenho de ultrassom multimodal do que modalidades isoladas na classificação de fibrose precoce\~\cite{Qin2024Fibrosis}. Esses trabalhos delimitam o escopo clínico: sem histopatologia pareada, a presente pesquisa mede ecogenicidade, não fibrose.

\section{Evolução Metodológica da Pesquisa}

A primeira fase concentrou-se na segmentação externa e na comparação entre arquiteturas. A limitação de dados motivou uma segunda fase de engenharia de fontes externas. Em seguida, uma estratégia inicial de pseudorrotulação utilizou limiares de confiança e filtros geométricos para selecionar máscaras. A inspeção dos casos mostrou, porém, que alta confiança não garante correção anatômica. A política foi modificada: os filtros passaram a priorizar revisão.

A quarta fase introduziu consenso entre U-Net e DeepLabV3. A concordância foi calibrada em esquema \textit{out-of-fold}, evitando medir estabilidade apenas em dados vistos. A quinta fase incorporou Medulla e, posteriormente, uma formulação multiclasse de córtex, medula e CEC. Uma expansão controlada de pseudorrótulos de Medulla foi então avaliada. Como o aumento do conjunto não melhorou o teste manual, a curadoria humana ganhou papel central.

A etapa seguinte transformou a curadoria em componente do sistema, com interface, banco SQLite, histórico de correções, status por camada e exportação. Finalmente, a pesquisa passou a medir ecogenicidade e revisou a narrativa sobre fibrose, separando marcador de imagem e desfecho clínico.

\section{Materiais e Métodos}

\subsection{Construção dos conjuntos supervisionados}

As anotações poligonais do kidneyUS são rasterizadas na resolução original. Para a cápsula, 486 registros do primeiro anotador correspondem a 468 imagens únicas. Após deduplicação por hash e agrupamento por exame, foram obtidas 328 imagens de treino, 70 de validação e 70 de teste. Para a tarefa multiclasse intrarrenal, 335 imagens foram divididas em 235, 50 e 50. Para Medulla binária, o conjunto possui 336 imagens, divididas em 236, 50 e 50.

\begin{table}[H]
\centering
\caption{Conjuntos supervisionados principais.}
\label{tab:datasets}
\begin{tabular}{lrrrr}
\toprule
\textbf{Tarefa} & \textbf{Total} & \textbf{Treino} & \textbf{Val.} & \textbf{Teste} \\
\midrule
Cápsula deduplicada & 468 & 328 & 70 & 70 \\
Intrarrenal multiclasse & 335 & 235 & 50 & 50 \\
Medulla binária & 336 & 236 & 50 & 50 \\
\bottomrule
\end{tabular}
\end{table}

\subsection{Deduplicação e prevenção de vazamento}

Hashes de conteúdo são utilizados para localizar arquivos repetidos. O identificador do exame é usado como chave de agrupamento durante a divisão dos conjuntos. Quando múltiplas máscaras estão associadas ao mesmo conteúdo, uma representação medóide pode ser escolhida pela similaridade com as demais. Esse procedimento reduz vazamento entre treino e teste.

\subsection{Pré-processamento}

Os modelos de segmentação recebem imagens redimensionadas e podem usar CLAHE. A imagem original é preservada e utilizada nas medidas quantitativas. Essa separação é necessária porque CLAHE redistribui intensidades locais e alteraria a variável de ecogenicidade.

\subsection{Benchmark da cápsula}

U-Net, UNet++, DeepLabV3-ResNet50 e SegFormer-B0 foram avaliados no mesmo teste deduplicado. Dice, IoU, precisão, recall e FPS foram usados como métricas. A U-Net foi selecionada como modelo principal da cascata.

\subsection{Engenharia das imagens externas}

Metadados do MONAI/NVIDIA foram filtrados por termos renais e retroperitoneais. Estudos candidatos foram baixados em lotes, DICOMs elegíveis foram convertidos para PNG em escala de cinza, frames RGB foram descartados por padrão e sequências cine foram reduzidas a frames representativos. No recorte utilizado, 4.479 imagens externas foram submetidas aos modelos.

\subsection{Pseudomáscaras e consenso entre modelos}

A confiança média dos pixels positivos é registrada como
\begin{equation}
C=\frac{1}{|M|}\sum_{i\in M}p_i.
\end{equation}
O mesmo caso é processado pela DeepLabV3 e a concordância é calculada por
\begin{equation}
S(M_U,M_D)=\frac{2|M_U\cap M_D|}{|M_U|+|M_D|}.
\end{equation}
A calibração é realizada em cinco folds agrupados por exame. Os limiares operacionais 0,89 e 0,94 são usados para ordenar a fila de revisão, não para validar automaticamente a anatomia.

\subsection{Segmentação intrarrenal e concordância humana}

U-Net e DeepLabV3 multiclasse predizem fundo, córtex, medula e CEC dentro da região renal. A avaliação é calculada por classe. As duas séries de anotação do kidneyUS também são comparadas nas imagens em que ambos os anotadores marcaram a estrutura, fornecendo contexto para a dificuldade da tarefa.

\subsection{Pseudo-expansão de Medulla}

DeepLabV3 e MedullaROIUNet geraram máscaras independentes. Candidatos concordantes foram filtrados e 342 pseudomáscaras foram adicionadas ao treino manual de 236 imagens. Validação e teste permaneceram exclusivamente manuais.

\subsection{Estabilidade da cascata}

O modelo de Medulla é avaliado duas vezes: com a cápsula manual e com a cápsula produzida pelo segmentador renal. Essa comparação quantifica a propagação de erro entre as etapas.

\subsection{Curadoria human-in-the-loop}

A interface permite revisar rim, córtex, medula e CEC separadamente. Os estados incluem pendente, aceita, corrigir, rejeitada e indisponível. Correções poligonais geram nova máscara sem sobrescrever a proposta original. O SQLite mantém imagem, revisor, camada, operação, timestamps, histórico e seleção das ROIs de ecogenicidade. A correção manual tem precedência sobre a máscara automática.

\subsection{Inferência integrada}

Uma nova imagem pode seguir o fluxo imagem $\rightarrow$ cápsula $\rightarrow$ ROI renal $\rightarrow$ U-Net intrarrenal $\rightarrow$ córtex/medula/CEC $\rightarrow$ ecogenicidade. A saída é armazenada separadamente e não entra automaticamente no treinamento.

\subsection{Ecogenicidade global}

Sejam $\mu_C$, $\mu_M$ e $\mu_{CEC}$ as médias de intensidade na imagem original. São calculadas:
\begin{equation}
R_{C/CEC}=100\frac{\mu_C}{\mu_{CEC}},
\end{equation}
\begin{equation}
D_{CEC-C}=\mu_{CEC}-\mu_C,
\end{equation}
\begin{equation}
CN=100\frac{\mu_{CEC}-\mu_C}{\mu_{CEC}+\mu_C}.
\end{equation}
Também são armazenados mediana, IQR, cobertura e saturação.

\subsection{ROIs pareadas córtex--CEC}

As máscaras de córtex e CEC são erodidas com elemento 5x5. A faixa vertical é dividida em três estratos e o sistema busca até um par de discos por estrato, com mesma área e profundidade aproximada. Para cada par $i$, calcula-se
\begin{equation}
R_i=100\frac{\mu_{C,i}}{\mu_{CEC,i}}, \qquad D_i=\mu_{CEC,i}-\mu_{C,i}.
\end{equation}
A mediana de $R_i$ é utilizada como estimativa principal e o IQR descreve dispersão. O revisor pode reposicionar as ROIs desde que permaneçam dentro das regiões válidas.

\section{Resultados}

\subsection{Segmentação da cápsula}

\begin{table}[H]
\centering
\caption{Segmentação da cápsula no teste deduplicado do kidneyUS.}
\label{tab:capsule}
\begin{tabular}{lccccc}
\toprule
\textbf{Modelo} & \textbf{Dice} & \textbf{IoU} & \textbf{Precisão} & \textbf{Recall} & \textbf{FPS} \\
\midrule
U-Net & \textbf{0,9290} & \textbf{0,8722} & \textbf{0,9236} & \textbf{0,9440} & 18,11 \\
SegFormer-B0 & 0,9223 & 0,8598 & 0,9117 & 0,9426 & \textbf{31,34} \\
DeepLabV3 R50 & 0,9203 & 0,8606 & 0,9219 & 0,9338 & 22,26 \\
UNet++ & 0,9099 & 0,8419 & 0,8910 & 0,9419 & 16,57 \\
\bottomrule
\end{tabular}
\end{table}

\subsection{Consenso OOF e triagem externa}

O consenso médio U-Net--DeepLabV3 foi 0,9416, com mediana 0,9627 e correlação 0,840 com o menor Dice frente à referência. Nas 4.479 imagens externas, a U-Net produziu 3.534 pseudomáscaras. Destas, 1.529 apresentaram consenso alto, 699 ficaram na faixa intermediária, 1.204 tiveram consenso inferior a 0,89 e 537 foram segmentadas por apenas um dos modelos. Em 510 imagens nenhum dos dois produziu máscara. Os números descrevem cobertura e estabilidade, não acurácia externa.

\subsection{Segmentação intrarrenal}

\begin{table}[H]
\centering
\caption{Segmentação intrarrenal no teste manual de 50 imagens.}
\label{tab:intrarenal}
\begin{tabular}{lccccc}
\toprule
\textbf{Modelo} & \textbf{Córtex} & \textbf{Medula} & \textbf{CEC} & \textbf{Dice médio} & \textbf{IoU médio} \\
\midrule
U-Net & \textbf{0,7075} & \textbf{0,7237} & 0,8469 & \textbf{0,7594} & \textbf{0,6163} \\
DeepLabV3 R50 & 0,6822 & 0,7156 & \textbf{0,8497} & 0,7492 & 0,6045 \\
\bottomrule
\end{tabular}
\end{table}

\subsection{Concordância entre anotadores}

\begin{table}[H]
\centering
\caption{Dice médio entre anotadores do kidneyUS.}
\label{tab:interobserver}
\begin{tabular}{lcc}
\toprule
\textbf{Estrutura} & \textbf{Imagens comuns} & \textbf{Dice} \\
\midrule
Capsule & 486 & 0,9461 \\
Cortex & 323 & 0,5961 \\
Medulla & 325 & 0,6936 \\
CEC & 468 & 0,8688 \\
\bottomrule
\end{tabular}
\end{table}

\subsection{Medulla: baseline, cascata e expansão}

A heurística antiga de regiões escuras apresentou Dice médio de 0,3291 contra o anotador 1. O DeepLabV3 supervisionado de Medulla elevou o Dice de teste para 0,7528 e a MedullaROIUNet obteve 0,7326. Quando a ROI manual foi substituída pela máscara prevista do rim, o DeepLab variou de Dice global 0,7626 para 0,7610, diferença de 0,0016. A intensidade média em Medulla apresentou correlação 0,9508 entre referência manual e predição em cascata.

\begin{table}[H]
\centering
\caption{Efeito da expansão por pseudorrótulos de Medulla.}
\label{tab:medulla_expand}
\begin{tabular}{lrrrr}
\toprule
\textbf{Configuração} & \textbf{Treino} & \textbf{Val. Dice} & \textbf{Teste Dice} & \textbf{IoU} \\
\midrule
Manual & 236 & 0,7235 & \textbf{0,7528} & \textbf{0,6035} \\
Manual + 342 pseudo & 578 & \textbf{0,7424} & 0,7523 & 0,6030 \\
\bottomrule
\end{tabular}
\end{table}

\subsection{Resultado histórico da base expandida}

Um experimento anterior com DeepLabV3-ResNet50 no \texttt{dataset_geral} expandido obteve Dice médio de teste 0,9366. Esse resultado é mantido como histórico e não é comparado diretamente ao benchmark deduplicado do kidneyUS, pois os conjuntos e protocolos são diferentes.

\section{Discussão}

A cápsula é a etapa mais estável do pipeline. O Dice de 0,9290 e a concordância humana de 0,9461 indicam uma referência relativamente consistente. Em contraste, córtex e medula apresentam menor desempenho automático e menor concordância entre anotadores, mostrando que o problema intrarrenal envolve também variabilidade da própria referência.

A avaliação da cascata fornece evidência de que a primeira etapa pode sustentar a segunda: no DeepLab de Medulla, substituir a ROI manual pela prevista alterou o Dice global em apenas 0,0016. Isso não elimina falhas em casos difíceis, mas mostra viabilidade do encadeamento.

A pseudo-expansão de Medulla é igualmente informativa. O treino cresceu de 236 para 578 imagens sem ganho no teste manual. Esse resultado é compatível com a literatura que enfatiza qualidade de pseudorrótulos em vez de quantidade\~\cite{Su2024ReliablePseudo,Huang2025SemiPLC} e justifica a curadoria \textit{human-in-the-loop}.

A ecogenicidade conecta anatomia revisada e caracterização quantitativa. A comparação córtex--CEC é motivada por medidas relativas descritas na literatura\~\cite{Zhao2025Echogenicity}. Entretanto, o protocolo de três pares, erosão e raio adaptativo é uma decisão operacional desta pesquisa e necessita análise de sensibilidade.

A imagem original é usada para ecogenicidade porque CLAHE altera a distribuição de intensidade. Mesmo assim, ganho, TGC, profundidade e pós-processamento permanecem fatores de confusão\~\cite{Manley2001}.

Por fim, a pesquisa não equipara ecogenicidade a fibrose. Trabalhos que predizem IFTA utilizam biópsia como referência\~\cite{Athavale2021}, e abordagens multimodais associam diferentes informações ultrassonográficas a histopatologia\~\cite{Qin2024Fibrosis}. O presente projeto ainda não possui esse padrão-ouro.

\section{Limitações e Ameaças à Validade}

As principais limitações são: tamanho do conjunto manual para classes internas; variabilidade interobservador; ausência de referência manual externa no MONAI; dependência de parâmetros de aquisição para ecogenicidade; ausência de eGFR, laudo ou histopatologia pareados em escala suficiente; e necessidade de medir a reprodutibilidade e o custo da própria curadoria.

\section{Continuidade da Dissertação}

O próximo ciclo deve transformar a infraestrutura em um protocolo experimental fechado. Um lote piloto deve ser revisado independentemente por especialista e revisor treinado, medindo concordância, tempo de revisão e taxa de correção. Depois, um teste totalmente revisado deve ser congelado e três cenários comparados: treino manual original, treino com pseudorrótulos automáticos e treino com pseudorrótulos aceitos ou corrigidos por humano.

Para ecogenicidade, devem ser avaliados raio das ROIs, número de pares, erosão, diferença vertical e seleção automática versus ajustada. O desfecho inicial deve ser reprodutibilidade. Em uma fase clínica posterior, uma coorte institucional poderá incorporar creatinina, eGFR, laudo e, quando disponível, histopatologia.

\section{Conclusão}

Esta pesquisa evoluiu de um benchmark de segmentação para uma infraestrutura rastreável de análise renal. O kidneyUS fornece a base supervisionada; a U-Net segmenta a cápsula com Dice 0,9290; o consenso com DeepLabV3 organiza a revisão de imagens externas; a segmentação intrarrenal produz córtex, medula e CEC; e a curadoria humana mantém decisões e correções auditáveis.

Os experimentos mostram que a cascata pode ser estável e que mais pseudorrótulos não garantem melhor generalização. A contribuição mais recente integra anatomia revisada e ecogenicidade quantitativa utilizando a imagem B-mode original e ROIs pareadas. O trabalho estabelece, assim, uma base metodológica para futura validação clínica sem antecipar conclusões que os dados atuais ainda não sustentam.

\bibliographystyle{sbc}
\bibliography{sbc-template}

\end{document}
